\section{Introduction}
\label{sec:intro}

User-generated video reaches its viewers only after the hosting platform has re-encoded it. The
platform decides, from properties of the upload that the uploader controls only indirectly, which
codec, which rate-control target and which ingest filtering to apply, and which rendition each
viewer receives~\cite{meta2021facebook}. None of these decisions is documented, and none can be
differentiated through. As a consequence, uploading the best possible master is not the same as
delivering the best possible playback.

Creators have developed heuristics for this situation. The best known is the \emph{tier trick}:
upscaling a 1080p video to 1440p or 4K before upload so that the platform assigns a more efficient
codec and a larger bit budget. Whether anything beyond this heuristic is available to an uploader
is an open question. Encoder-side pre-processing is well studied: learned pre-processors in front of
standard encoders reduce bitrate at equal quality~\cite{bourtsoulatze2020deep,chadha2021deep,
ma2023rate,zhao2025tridynamic}. These methods, however, are deployed by the party that runs the
encoder and knows its settings. The uploader faces a different problem: the encoder, its settings
and the ladder that selects them are hidden, and every evaluation costs a real upload.

\paragraph{This paper.}
We formulate uploader-side pre-compensation as black-box optimisation of a transform
$\Gt$ applied before upload, so that the playback $\Y = \fp(\Gt(\X))$ is close to the source
$\X$ (\cref{sec:problem}). We study two routes: sample-efficient search over an interpretable
seven-parameter filter chain, and a learned residual network trained through a differentiable
emulator of the platform (\cref{sec:method}). Our evaluation is designed around two pitfalls that we
found to dominate the outcome:
\begin{enumerate}
  \item \textbf{Tier gain versus filter gain.} Most of the improvement comes from selecting a
        better encoding tier. We decompose every result into the tier gain and the filter gain on
        top of it.
  \item \textbf{Objective versus judge.} An optimiser finds whatever its objective rewards. We
        separate the optimised metric from independent judge metrics, and we count a gain only if
        the judges agree~\cite{zvezdakova2019hacking,siniukov2021hacking}.
\end{enumerate}

\paragraph{Contributions.}
\begin{itemize}
  \item A formulation of uploader-side pre-compensation against an adaptive, non-differentiable
        platform pipeline, with a decomposition of the gain into a tier component and a filter
        component (\cref{sec:problem}).
  \item Evidence, on simulated platforms, that the filter gain found by optimising a composite of
        VMAF NEG, LPIPS and CAMBI is not confirmed by SSIMULACRA2 and XPSNR, and that every
        single-metric objective degrades at least one judge it does not optimise (a ``Goodhart
        matrix'', \cref{sec:results-p0}).
  \item A documented failure mode of emulator-based training (exploitation of the emulator) and its
        mitigation by a trust region (\cref{sec:results-p0}).
  \item A pre-registered, content-scoped test of a learned pre-compensation with a content gate,
        evaluated once on held-out confirmation clips (\cref{sec:results-p1}).
        \pending{state the outcome.}
  \item \pending{Real-platform measurements: ladder probe, closed loop and viewer study
        (\cref{sec:results-p2}).}
  \item An open-source implementation, clip manifests with per-clip licenses, and the
        pre-registration files.
\end{itemize}
